\section{Witness vectors and mixing matrices}\label{sec:moments}

A witness must be distinguishable from the directions whose images we
have fixed. To see how to choose it, regard a vector of
\(B=\mathsf S\otimes F\) as a matrix with rows indexed by selector
coordinates and columns indexed by base coordinates. A single point
\(p_s\otimes\bar z\) has rank 1 as such a matrix. We will also encounter
a different rank, that of a symmetric matrix acting on \(B\) itself.
The next two lemmas relate the two descriptions.

Write \(\mathsf F=F\), and let \(W\) be the span of the point matrices
on \(B\), allowing every tag. The linear forms \(\epsilon_{\mathsf S}\)
and \(\epsilon_{\mathsf F}\) read the constant coordinate in each factor.

\subsection{The span of the point matrices}

% lax begin Lax793645.AffinePointMatrices
\begin{lemma}\label{lem:base-moments}
The span of $(1,z)(1,z)\trans$ consists precisely of symmetric matrices
$Z$ satisfying $\operatorname{diag}Z=Z\epsilon_{\mathsf F}$.
Consequently $W$ is the tensor product of the corresponding affine
point-matrix spaces on $\mathsf S$ and $\mathsf F$. Its tensors are
separately symmetric in the selector and base indices and satisfy the
Boolean diagonal identity in each factor.
\end{lemma}
% lax end
% lax begin Lax793645Proofs.Manuscript.AffinePointMatrices.mem_span_pointSquare_iff
\begin{proof}
Write \(P(z)=(1,z)(1,z)\trans\). Every \(P(z)\) is symmetric and has
\(P(z)_{ii}=P(z)_{i0}\). Conversely, let \(Z\) have these properties.
Choose an ordinary coordinate \(i\). Set \(u_i=1\) and
\(u_j=Z_{ji}\) for the other ordinary coordinates. The matrix
\[
 Z-P(u)-(Z_{ii}-1)P(e_i)
\]
has zero row and column at \(i\). It still satisfies the same identities,
and previously cleared rows and columns remain zero. Repeat at the remaining
coordinates. The matrix left at the constant coordinate is a multiple of
\(P(0)\). This proves the span assertion.
% lax begin Lax793645Proofs.Manuscript.AffinePointMatrices.affineTensor_identities
The tensor-product assertion follows because the selector
and base points vary independently. The identities in another basis
use the coordinates of $\epsilon_{\mathsf S},\epsilon_{\mathsf F}$.
% lax end
\end{proof}
% lax end

% lax begin Lax793645.SelectorConfinement
\begin{lemma}\label{lem:selector-confinement}
If $w$ is a symmetric matrix on $\mathsf S\otimes\mathsf F$, separately
symmetric in its two pairs of indices, and $\rank w\le r$, then
\[
 \operatorname{col}(w)\subseteq U\otimes\mathsf F
 \quad\text{for some }U\subseteq\mathsf S,\quad \dim U\le r^2.
\]
Consequently each vector in the column space has rank at most \(r^2\)
when reshaped as a matrix with one row per selector coordinate.
\end{lemma}
% lax end
% lax begin Lax793645Proofs.Manuscript.SelectorConfinement.selector_confinement
\begin{proof}
For a column of the tensor matrix indexed by $(b,j)$, reshape it as the
matrix $C_{ai}=w_{(a,i),(b,j)}$. Separate symmetry gives
$C_{ai}=w_{(a,j),(b,i)}$, so $C$ is a submatrix of $w$ and has rank
at most $r$. Choose at most $r$ columns spanning the original column space.
The sum of their reshaped column spaces has dimension at most $r^2$
and contains every selector slice.
\end{proof}
% lax end


\subsection{Protected subspaces}

Fix the pin budget \(K\) and an upper bound \(t_* =28\) for the number of
selected atoms. Choose
\[
 R\ge4(K+t_*),\qquad q_*=R^2+t_*.
\]
The extra room in \(R\) covers the joint obstruction at both endpoints.
In each primal pin space consider the span of the vectors whose rank
after reshaping is at most \(q_*\). Choose a basis of this span from
among those vectors. Let \(S_0\) be the sum of the column spaces of
the reshaped basis vectors. For \(F_0\), take the sum of the row spaces
of a basis of each \emph{whole} pin space. Take these sums over the
components of one endpoint, so the same \(S_0,F_0\) serve all its
components. If \(P_0\) bounds the sum of the pin dimensions, then

\begin{equation}\label{eq:label-budget}
 \dim S_0\le q_*P_0,\qquad \dim F_0\le(b+1)P_0.
\end{equation}
Set
\[
 D_0=S_0\otimes F_0,\qquad
 L=S_0\otimes\mathsf F+\mathsf S\otimes F_0.
\]
Every pin vector lies in \(L\). Every pin vector with reshaped rank at
most \(q_*\) lies in \(D_0\). The first bound in
\eqref{eq:label-budget} is independent of \(b\), and the second is independent
of \(n\).

% lax begin Lax793645.SelectedAtoms
\begin{lemma}\label{lem:protected-atoms}
Let $k_i=p_i\otimes v_i$, $1\le i\le t$, where
$\epsilon_{\mathsf S}(p_i)=\epsilon_{\mathsf F}(v_i)=1$.
Suppose the $p_i$ are independent modulo $S_0$ and the $v_i$ are
independent modulo $F_0$. Put $A=\operatorname{span}\{k_i\}$.
Then $A\cap L=0$. If $w\in W$ and
\[
 w\in(D_0+A)\otimes B+B\otimes(D_0+A),
\]
then
\[
 w=\sum_i c_i k_i k_i\trans+w_0,\qquad w_0\in L\otimes L.
\]
\end{lemma}
% lax end
% lax begin Lax793645Proofs.Manuscript.SelectedAtoms.tensor_atom_isolation
\begin{proof}
Extend bases of $S_0,F_0$ by the selected vectors and then to full bases.
The selected simple tensors occupy distinct diagonal positions outside
the rows and columns belonging to $L$, proving $A\cap L=0$.
An entry of $w$ must have an index in $D_0$ or at a selected position.
Separate symmetry moves the entry between $p_i\otimes v_i$ and
$p_a\otimes v_b$ to that between $p_i\otimes v_b$ and $p_a\otimes v_i$.
It is zero unless $a=i$ or $b=i$. If $a=i,b\ne i$, the diagonal entry
at $p_i\otimes v_b$ is zero. Its base diagonal identity has just one
potentially nonzero term, with coefficient
$\epsilon_{\mathsf F}(v_i)=1$, so that term is zero too.
The selector identity treats the other case. Each selected row and
column therefore has only its diagonal entry. Subtract these point
matrices. Every remaining entry has a $D_0$ index. If its other index
were outside $L$, separate symmetry would put both indices outside
$D_0+A$, forcing zero. Thus the residual tensor lies in $L\otimes L$.
\end{proof}
% lax end


A span of dimension \(d\) contains at most \(2^d\) affine selector vectors.
Thus the exclusions needed to choose the selectors independently modulo
\(S_0\) have a bound fixed before \(b\). Include the other finite label
exclusions used below in a budget \(B_{\rm exc}\), and choose
\begin{equation}\label{eq:selector-size}
 B_{\rm exc}<2^{b-2},\qquad 3B_{\rm exc}<2^b.
\end{equation}
When several lists are chosen, the exclusion budget includes all previously
chosen selectors. It is deliberately larger than needed for one list.

Base independence uses the free block \(Z\). Conditional on earlier point
parameters, a subspace of dimension at most \(\dim F_0+t_*\) contains at
most \(2^{\dim F_0+t_*}\) of the possible \(Z\)-vectors. Hence, at every
orientation, failure of independence modulo \(F_0\) has probability at most
\begin{equation}\label{eq:base-rejection}
 \delta_n=t_*2^{\dim F_0+t_*-n}=2^{-n+O(1)}.
\end{equation}
This probability is conditional on the orientation, and the bound is
the same for every orientation. Discarding these inputs therefore costs
at most \(\delta_n\) under any mixture of orientations.

\subsection{The input laws}

At tag \(l\), a generic query draws all its allowed coordinates
uniformly and independently. A shared point is obtained by setting its
ordinary coordinates to zero. We use shared points when evaluating
identities that hold for every generic point. Occasionally we prescribe
the few coordinates on which specified testers depend. Their number is
fixed, so the event has probability bounded below independently of \(n\).

The diagonal bit of an atom is \(q=\eta(vv\trans)\). Its role bit is
\(a(A)\).

Together with these bits, a query can record the individual
tester values and the evaluations \(T(A,c)\) on a basis of an effective
space. Each of these unary bits is a quadratic function of the point
parameters of bounded rank. The same is true of \(u_{o'}U_oA\), since
there are only \(h\) channel products per component. Expanding indicators
of finitely many such bits gives a bounded Fourier \(\ell^1\)-norm,
uniformly in \(n\) and in the orientation.

\subsection{Unary mixing}

% lax begin Lax793645.MixingMatrices
\begin{lemma}\label{lem:generic-mixers}
Choose a sufficiently large fixed \(J\). For all large \(n\), the reduced
family in \eqref{eq:reduced-mixer} can be chosen so that the following holds.
For every nonzero profile \(x\) of total rank at most \(K\), every tag and
label, the polar rank in the \(Z\) variables of \(T(x,A(l,s,z))\)
is at least \(2(J-s_0)\). Here \(s_0\) is fixed before \(J\).
The expression uses a bounded-dimensional space of linear forms on \(Z\).
The conclusion survives every fixing of the other coordinates.
\end{lemma}
% lax end
% lax begin Lax793645Proofs.Manuscript.MixingMatrices.goodReducedMixers
\begin{proof}
A symmetric rank-\(r\) matrix has a factorization \(ZHZ\trans\), with
\(Z\) of column rank \(r\) and \(H\) nonsingular and symmetric.
Indeed, choose a column basis \(Z\), a left inverse \(L\), and take
\(H=LxL\trans\). Thus profiles of total rank at most \(K\) number
at most \(2^{A_0n}\), for a fixed \(A_0\).

Choose the entries involving \(Z\) in all \(L_j^c,R_j^c\) at random.
For a tag \(l\), the forward profile coefficients are \(x_{e_c}\) for
\(c\ne l\). The reverse coefficients are \(z_c(x)\) for edges
\(e_c\) incident with \(l\). These coefficient families determine a
nonzero cut profile. In fact a nonzero cut has at least 2 nonzero cycle
components, since their sum is zero. Some nonzero \(x_{e_c}\) has \(c\ne l\).
Factor the nonzero coefficients as above. Treat the resulting linear
forms in each \(L\) and \(R\) as formal variables. For each \(j\) a
nonzero bilinear product remains, so their formal quadratic form has polar
rank at least \(2J\).

The list of formal variables has bounded length \(m_0\). Forward and
reverse restrictions of the same random matrix agree where their
rectangular blocks intersect. There are only boundedly many such
compatibility bits. Consequently their distribution is bounded by
\(2^{C_0}\) times the uniform law on all \(m_0\)-by-\(n\) matrices.
The constants \(m_0,C_0\) depend on the fixed parameters, not on \(n\).
Failure to obtain the desired rank costs at least \(s_0+1\) linear
relations among the rows. There are at most \(2^{m_0(s_0+1)}\) choices
of independent relations. For a uniform matrix each choice has
probability \(2^{-(s_0+1)n}\), giving the bound
\[
 2^{m_0(s_0+1)-(s_0+1)n}.
\]
Choose \(s_0>A_0\) and sum this bound over all profiles, tags, and labels.
The probability tends to zero. Restricting a bilinear form to a subspace
of codimension \(s_0\) loses at most \(2s_0\) rank. This proves the
claim. Other input coordinates change only affine terms, and \(T^0\)
has no \(Z\) variables.
\end{proof}
% lax end

A quadratic form of polar rank \(r\) has Walsh bias at most \(2^{-r/2}\).
Fourier inversion on the bounded list of effective-basis bits therefore
gives a fixed positive probability for each consistent basis assignment,
after the role and tester data are fixed. We choose \(J\) large enough
for this lower bound to be uniform over all effective spaces of dimension
at most \(K^2\). This is the unary positivity estimate used later.

\subsection{Explicit binary mixing}

Enumerate the \(L\) and \(R\) matrices by \(F_0,\ldots,F_{r-1}\), where
\(r=2Jg\). Split \(\#\) into two parts of size \(d=\lfloor n/2\rfloor\),
and let \(S\) be the nilpotent shift on \(\F^d\). In the constant-selector
\(\#\)-block of \(F_i\), put
\[
 \begin{pmatrix}0&S^i\\0&0\end{pmatrix}.
\]
Set its other selector \(\#\)-blocks to zero. For \(0\le\rho<b\), use
the same matrix with exponent \(r+\rho\) as \(M_\rho\), and define
\[
 v(s,z)\trans E^\#v(t,z')
   =\sum_\rho(s_\rho+t_\rho)z_\#\trans M_\rho z'_\#.
\]
This vanishes at a single point and leaves the tester identity intact.
No entry involving \(Z\) was changed, so unary mixing is preserved.

For distinct labels, any nonzero required parity of the ordered binary
evaluations has a \(\#\)-restriction of the form
\[
 \begin{pmatrix}0&P(S)\\Q(S)\trans&0\end{pmatrix},
\]
with at least one of \(P,Q\) nonzero. The exponent ranges for the indexed
mixers and for \(E^\#\) are disjoint. The latter coefficients involve
\(s_\rho+t_\rho\), not all zero. A nonzero polynomial of degree at most
\(r+b-1\) satisfies
\[
 \rank P(S)\ge d-(r+b-1).
\]
Indeed, if \(S^a\) is the lowest power that occurs, write
\(P(S)=S^a(I+ST(S))\). The second factor is invertible because \(S\)
is nilpotent, and \(S^a\) has rank \(d-a\). For large \(n\) the displayed
rank is at least \(n/5\). This proves binary mixing with exponentially
small error.

An invertible change of the \(Z\) coordinates at one vertex also
preserves the raw frame law. Indeed, \(E\) has no \(Z\) entries, so
precomposing the frame with this change leaves its defining equations
unchanged. Changing the point parameters correspondingly preserves the
keys. We will use this symmetry when some linear functions of the keys
have already been fixed.
